\section{Selecting the two auxiliary primes}\label{sec:hosts}

Fourier projection leaves terms twisted by either of two auxiliary quadratic
characters or by their product. Their untwisted factors are the primitive characters
inducing powers of the fixed primitive character \(\psi\).
We select the auxiliary primes in prescribed progressions
using correlation estimates for the prime sums and a
simultaneous exclusion of bad pairs. For the selected pair
the unwanted modes are small throughout the interval used
in Section~\ref{sec:comparison}. Appendix~\ref{sec:cohomology}
then supplies a cohomology subinterval for this same pair. Definition~\ref{hyp:host-residue-class} specifies the congruence and size conditions on the candidate primes. The pair is selected from the two logarithmic intervals in \eqref{eq:host-pair-log-windows}.

\subsection{Correlation in a progression}

Fix the progression from which a prime will be chosen. We first average the Jacobi symbol over its candidate primes for each admissible squarefree denominator, allowing one exceptional denominator. The choice of a pair for the whole interval comes later.

\begin{paramdefinition}[Parameter range for the progression average]\label{hyp:host-residue-class}
Fix \(0<\varepsilon<1/52\) as in Proposition~\ref{prop:quadratic-prime-pair}. Let \(Y_0\) be the prime selection scale, \(J\ge1\) a fixed exponent,
and \(Z=Y_0^{400J}\). Let \(m_0,a_0\) be integers with \(8\mid m_0\) and \(\gcd(a_0,m_0)=1\). The primes under consideration form the set
\[
 \mathcal P=\{l\text{ prime}:\sqrt U<l\le U,\ l\equiv a_0\pmod{m_0}\}.
\] The admissible range of \(U\) for this modulus is
\[
 \log U\ge\tfrac14(\log Y_0)^2\log m_0,\qquad
 Y_0^{\varepsilon/6}\le\log U
       \le Y_0^{\varepsilon/5}(\log Y_0)^2 .
\]
The moduli in this range are restricted to those for which each prime divisor of \(m_0\) is at most
\(Y_0^{\varepsilon/5}\) or exceeds \(Z\).
\end{paramdefinition}
For the first selection take \(m_0=8\) and
\(\log U=Y_0^{\varepsilon/5}\); all inequalities hold for large
\(Y_0\). For the second selection \eqref{eq:host-pair-log-windows}
and the modulus calculation following it verify the same range
with the first prime already fixed.

\begin{proposition}[Quadratic correlation in one progression]\label{prop:quadratic-host-correlation}
For the data of Definition~\ref{hyp:host-residue-class}, for any fixed \(B>0\)
and all sufficiently large \(Y_0\), the relative
prime number theorem and the exceptional-character alternative give
\begin{equation}\label{eq:quadratic-corr}
       \left|\sum_{l\in\mathcal P}(l/s)\right|
       \le Y_0^{-B}\abs{\mathcal P}
\end{equation}
for all odd squarefree \(1<s\le Z^2\) coprime to \(m_0\), with at most
one exceptional value of \(s\). Here \((\cdot/s)\) is the Jacobi symbol.

\end{proposition}

\begin{proof}
\proofstep{Use the relative progression estimate.}

To obtain \eqref{eq:quadratic-corr}, use the relative progression
estimate at modulus \(m_0s\). The displayed hypotheses imply
\(\log(m_0s)=o(\log U)\) uniformly for \(s\le Z^2\).
For example, Corollary~1.4 of Thorner and Zaman applies at all
intermediate scales \(\sqrt U\le t\le U\), since
\(t\ge(m_0s)^{12}\) for sufficiently large \(Y_0\).

To pass from logarithmic weights to a count of primes, keep the exceptional character in the main term. Put \(m=m_0s\), \(V=\sqrt U\), and, for a reduced class \(b\bmod m\), set
\[
 h_b(t)=1-\xi(b)t^{\beta-1},\qquad F_b(t)=\int_0^t h_b(v)\,dv .
\]
If the modulus has no exceptional character use \(h_b=1\). For \(V\le t\le U\), Corollary~1.4 of Thorner and Zaman gives
\[
 \sum_{p\le t,\ p\equiv b\ (m)}\log p
     =\frac{F_b(t)}{\varphi(m)}(1+O(E)),\qquad
 E=e^{-c\log U/\log m}+e^{-c\sqrt{\log U}}
\]
with an absolute \(c>0\), decreased to allow the lower endpoint.
For \(\xi(b)=1\), the function \(h_b\) increases and
\(0\le F_b(t)\le t h_b(t)\) in this range. The lower inequality follows from
\(t^{1-\beta}\ge1/\beta\) for \(\log t\ge2\), since an exceptional zero has \(\beta>1/2\).
For \(\xi(b)=-1\), \(F_b(t)\le3t\) and \(h_b(t)\ge1\).

Thus partial summation gives
\[
 \#\{V<p\le U:p\equiv b\ (m)\}
 =\frac{I_b}{\varphi(m)}(1+O(E)),\qquad
 I_b=\int_V^U h_b(t)\frac{dt}{\log t}.
\]
Indeed its two endpoint errors and integral error together are
\(O(EU h_b(U)/(\varphi(m)\log U))\). This is a relative error:
the last half of the interval supplies
\[
 I_b\ge \int_{U/2}^U h_b(t)\frac{dt}{\log U}
       \ge \frac{U h_b(U)}{4\log U}.
\]
For the plus character value this follows from
\(1-e^{-z/2}\ge(1-e^{-z})/2\) with \(z=(1-\beta)\log U\).
\cite[Corollary~1.4 and (1.7)]{thorner-zaman}

Separate the two integrals in the exceptional case:
\[
 I=\int_V^U\frac{dt}{\log t},\quad I_\beta=\int_V^U t^{\beta-1}\frac{dt}{\log t},
 \qquad I_b=I-\xi(b)I_\beta .
\]

Under the Chinese remainder identification write
\(b=(a_0,c)\), \(c\in(\Z/s)^\times\), and
\(\xi(b)=\xi_0(a_0)\xi_s(c)\). The ordinary weighted
main term contains \(\sum_c(c/s)=0\). The exceptional
weighted term contains
\[
   \xi_0(a_0)\sum_{c\in(\Z/s)^\times}(c/s)\xi_s(c),
\]
which vanishes unless \(\xi_s=(\cdot/s)\). In that case
\(\sum_c\xi_s(c)=0\), so the unweighted main term is \(I/\varphi(m_0)\). Since \(t\ge\sqrt U\) throughout the integral,
\[
  0\le I_\beta/I\le U^{-(1-\beta)/2}.
\]
In both cases the sum of the errors is bounded by the
relative error factor times the unweighted main term,
because the individual progression main terms are nonnegative.

\proofstep{Use a single exceptional-character range.}

For all moduli take \(D=m_0Y_0^{800J}\). The lower bound on
\(\log U\) gives the two estimates
\[
 \log m_0\le\frac{4\log U}{(\log Y_0)^2},\qquad
 800J\log Y_0=o_J\!\left(\frac{\log U}{(\log Y_0)^2}\right).
\]
Thus, for large \(Y_0\),
\[
       \frac{\log U}{\log D}\ge\tfrac15(\log Y_0)^2 .
\]
If \(1-\beta\ge c/\log D\), the exceptional contribution therefore satisfies
\[
 I_\beta/I\le\exp\!\left(-\frac{c\log U}{2\log D}\right)
     \le \exp\!\left(-\tfrac c{10}(\log Y_0)^2\right)\ll_B Y_0^{-B}.
\]
Here \(B\) is any fixed exponent.

For nearer zeros, the Landau--Page alternative gives the same
primitive real character for all moduli up to \(D\).
Its restriction can equal the primitive quadratic character of
at most one such \(s\). Indeed, the restrictions for two squarefree
values differ on a prime dividing exactly one of them, as is seen
by taking a nonsquare there and the unit residue at the other primes.
Finally apply the same interval estimate at \(m_0\).
In the worst class the lower bound for \(I_b\) on \([U/2,U]\) gives
\[
 |\mathcal P|\gg\frac{U}{\varphi(m_0)\log U}
                 \min\{1,(1-\beta)\log U\}.
\]
Without an exceptional zero the last factor is one. With one, use
\(1-\beta\gg m_0^{-1/2}\) from \cite[(1.2), with exponent \(1/2\)]{thorner-zaman}.
Since \(\log m_0=o(\log U)\), this lower bound exceeds \(U^{9/10}\) for all sufficiently large \(Y_0\). These uses of relative prime
distribution and Page uniqueness concern the whole range of moduli.
\hfill\mbox{\cite{page,thorner-zaman}}\label{cor:cohomology-host-windows-1}

\end{proof}

\subsection{A mean square over the whole scale interval}

A bound at one scale would allow the exceptional primes to vary with the scale. To keep one set of primes throughout, we estimate the variation of the truncated sum.
For coefficients \(\abs{\kappa(n)}\le1\) put
\[
 S_l(y)=\sum_{n\le Z}\Lambda(n)\min(1,y/n)\,\kappa(n)(n/l),
 \qquad Y_0\le y\le Y_0^J .
\]
\begin{proposition}[Uniformity over the scale interval]\label{prop:quadratic-host-uniform}
For the data of Definition~\ref{hyp:host-residue-class}, fix \(\tau>0\) and a
list of bounded cardinality of coefficient sequences with
\(|\kappa(n)|\le1\) for every \(n\).
Then all but an \(O_\tau(Y_0^{-\tau})\)
fraction of \(\mathcal P\) satisfy, simultaneously for every sequence in the list, 
\[
 \lvert S_l(y)\rvert\le y^{1/2+\tau}/2\qquad(Y_0\le y\le Y_0^J).
\]

\end{proposition}

\begin{proof}
\proofstep{Truncation and coefficient bounds.}

Write \(A(t)=\sum_{j\ge2,\,p^j\le t}\log p\). Chebyshev's estimate
\(\sum_{p\le v}\log p\ll v\) gives, for \(t\ge2\),
\[
 A(t)\ll \sqrt t+t^{1/3}\log t\ll\sqrt t .
\]
Indeed the term with \(j=2\) gives the first bound and there
are at most \(\log_2 t\) remaining powers to sum. Including the
weighted tail, the unsigned contribution of higher powers is at most
\[
 A(y)+y\int_{(y,Z]}t^{-1}\,dA(t)
   =yA(Z)/Z+y\int_y^Z A(t)t^{-2}\,dt\ll\sqrt y .
\]
The prime contribution from \(p\mid m_0\) is
\(O(Y_0^{\varepsilon/5})\), since the large divisors exceed
\(Z\). For the other primes write
\(a_p=(\log p)\min(1,y/p)\kappa(p)\).
Chebyshev's bound and partial summation give
\[
 \sum_p|a_p|^2\ll
  \sum_{p\le y}(\log p)^2+y^2\sum_{p>y}\frac{(\log p)^2}{p^2}
  \ll y\log y,\qquad
 \sum_p|a_p|\ll y\log Z .
\]

\proofstep{Diagonal, exceptional pair and remaining pairs.}

On expanding \(\sum_l\abs{S_l(y)}^2\), the diagonal is bounded by
\(\abs{\mathcal P}\sum\abs{a_p}^2\). For distinct primes \(p,r\)
the product of symbols is \((l/(pr))\) times a fixed sign by
quadratic reciprocity, since the class modulo \(8\) is fixed.
The exceptional value in \eqref{eq:quadratic-corr} can correspond
to at most one unordered pair \(\{p,r\}\), which costs at most
\(\abs{\mathcal P}(\abs{a_p}^2+\abs{a_r}^2)\).
The other pairs cost \(\ll Y_0^{-B}\abs{\mathcal P}y^2(\log Z)^2\).
With \(B>J\), these estimates give
\[
 \sum_{l\in\mathcal P}\abs{S_l(y)}^2\ll\abs{\mathcal P}\,y\log y .
\]
\proofstep{Bounding the proportion of primes exceptional somewhere in a dyadic interval.}

This bound is uniform in the coefficients. Since \(yS_l'(y)\)
is the same sum restricted to \(n>y\) almost everywhere,
\(\sum_l\abs{S_l'(y)}^2\ll\abs{\mathcal P}\log y/y\).
For \(V\le W\le2V\) in the interval and arbitrary
\(w_l\in[V,W]\), the finite sum \(S_l\) is absolutely
continuous. Hence \(S_l(w_l)-S_l(V)=\int_V^{w_l}S_l'(y)\,dy\).
The square of this increment is bounded by
\((W-V)\int_V^W|S_l'(y)|^2\,dy\). Summing over \(l\)
gives
\[
 \sum_l\abs{S_l(w_l)}^2
   \le 2\sum_l\abs{S_l(V)}^2+
       2V\int_V^W\sum_l\abs{S'_l(y)}^2\,dy
   \ll\abs{\mathcal P}\,V\log(2V).
\]

Fix \(\tau>0\). For primes at which the bound
\(\abs{S_l(y)}\le y^{1/2+\tau}/2\) fails anywhere in
\([V,W]\), choose \(w_l\in[V,W]\) with \(\abs{S_l(w_l)}>w_l^{1/2+\tau}/2\). The fraction of these
primes is \(O(V^{-2\tau}\log(2V))\). Summing across dyadic blocks
shows that all but \(O_\tau(Y_0^{-\tau})\) of the primes meet the
bound for every \(y\in[Y_0,Y_0^J]\).
A union bound applies to any list of bounded cardinality.

\end{proof}

\subsection{The full smoothed sum and the shared interval}

The mean square controls a truncated polynomial. The Fourier comparison uses the full
smoothed sum and several twists at once. We now impose these requirements on the same
primes.

\begin{auxlemma}[Logarithmic derivative inside a zero-free rectangle]\label{lem:host-log-derivative}
Let \(Q\ge4\), \(3/4\le\alpha<1\), and \(0<\nu<1-\alpha\).
Suppose \(\chi\) is primitive and nonprincipal, of conductor \(f\le Q\).
If \(L(s,\chi)\) has no zero in
\(\RePart s\ge\alpha,\ |\Ima s|\le Q\), then
\[
 \left|\frac{L'}{L}(\sigma+it,\chi)\right|
   \ll(1+\nu^{-1})\log Q
     \qquad(\sigma\ge\alpha+\nu,\ |t|\le Q-1)
\]
with an absolute implied constant.
\end{auxlemma}

\begin{proof}
For \(\sigma\ge2\) the absolute Euler series gives the bound.
For \(\alpha+\nu\le\sigma\le2\) write \(s=\sigma+it\),
\(s_0=2+it\) and \(\rho=\beta+i\gamma\) for a nontrivial zero,
counted with multiplicity. Subtract the logarithmic derivatives of
the Hadamard product at \(s\) and \(s_0\). The conductor factor cancels.
The gamma factors contribute \(O(1)\), by Stirling's formula for
\(|t|\ge2\) and compactness for \(|t|\le2\). Consequently
\[
 \frac{L'}{L}(s,\chi)=O(1)+\sum_\rho
        \left(\frac1{s-\rho}-\frac1{s_0-\rho}\right).
\]
The unit-interval zero count gives
\(\#\{\rho:j\le|\gamma-t|<j+1\}\ll\log(f(3+|t|+j))\)
for integers \(j\ge0\). Both this estimate and the product are
given in \cite[Lemma~10.3 and Section~10.2]{classical-region}.

For \(j=0\) use \(|s-\rho|\ge\nu\) and \(|s_0-\rho|\ge1\).
For \(j\ge1\) use the difference of the two fractions. Its
absolute value is at most \(2/j^2\). The two contributions are
bounded, respectively, by
\[
 (1+\nu^{-1})\log(f(3+|t|)),\qquad
 \sum_{j\ge1}\frac{\log(f(3+|t|))+\log(1+j)}{j^2}
    \ll\log(f(3+|t|)).
\]
Now use \(f\le Q\) and \(|t|\le Q-1\).
\end{proof}

\begin{proposition}[Removing the truncation for a bounded list of twists]\label{prop:quadratic-host-smoothed}
Assume Definition~\ref{hyp:host-residue-class}. Fix \(\tau>0\)
and an upper bound for the cardinality of the list of primitive base
characters. Their conductors may vary subject to \(f\le F\le\sqrt U\).
There is a subset \(\mathcal P'\subseteq\mathcal P\) with
\(|\mathcal P\setminus\mathcal P'|/|\mathcal P|=o(1)\) as
\(Y_0\to\infty\), uniformly over such lists. For every \(l\in\mathcal P'\)
and each twist \(\chi\) of a listed character by \((\cdot/l)\),
\[
 \lvert G_\chi(y)\rvert\le y^{1/2+\tau}\qquad(Y_0\le y\le Y_0^J).
\]

\end{proposition}

\begin{proof}
\proofstep{Zero density is applied to primitive twists.}

For a primitive base character of conductor \(f\le F\le\sqrt U\),
the twist at \(l\in\mathcal P\) is primitive of conductor \(fl\).
Put \(Q=FU\) and \(\alpha=1-1/400\). Jutila's estimate (1.6) gives
\[
 \sum_{q\le Q}\ \sum_{\chi\bmod q}^{*}N(\alpha,Q,\chi)
     \ll (Q^3)^{2(1-\alpha)+1/100}\ll Q^{1/4} .
\]
The star denotes primitive characters and \(N(\alpha,T,\chi)\)
counts zeros with real part at least \(\alpha\) and height at most
\(T\) in absolute value. The exponent in Jutila's bound uses
\(\alpha\ge21/26\) and \(Q\ge2\), with the parameter \(k=2\).
Hence there are \(O(Q^{1/4})\) primitive characters of conductor
at most \(Q\) with a zero in that rectangle. The conductor \(fl\)
determines \(l\), as \(l>F\). Consequently the excluded primes
are \(O(U^{3/8})\), a vanishing fraction of \(\mathcal P\).
\hfill\mbox{\cite{jutila}}

\proofstep{Estimate the tail by a contour difference.}

For a remaining twist \(\chi\), write
\[
  R_\chi(z)=-\ell_\chi(1)-\sum_{n\le z}\Lambda(n)\chi(n)/n .
\]
Here \(\ell_\chi=L'/L(\,\cdot\,,\chi)\). Since \(f<l\) and the
quadratic twist is nontrivial at \(l\), the primitive twist is
nonprincipal. Lemma~\ref{lem:host-log-derivative} applies with
\(\alpha=1-1/400\) and \(0<\nu\le1/800\). Use the contour formula for \(S_\chi(z)\) on
\(\RePart s=2\) and take the difference at \(z+h\) and \(z\),
where \(0<h\le z\). Shift to \(1-1/400+\nu\) at heights
\(\pm Q/2\). The residue at \(1\) is \(-h\ell_\chi(1)\).

On the new edge \(\RePart s=\sigma_1=1-1/400+\nu\),
\[
 (z+h)^s-z^s=s\int_z^{z+h}u^{s-1}\,du,\qquad
 \int_{-Q/2}^{Q/2}\frac{dt}{|\sigma_1+it-1|}
      \ll_\nu\log Q .
\]
Since \(0<h\le z\), the first integral cancels the factor
\(s\) and leaves size at most \(h z^{\sigma_1-1}\).
The horizontal edges and the two tails on \(\RePart s=2\)
give \(O_\nu((\log Q) z^2/Q)\). Thus for
\(\sigma_1=1-1/400+\nu\)
\[
 \abs{S_\chi(z+h)-S_\chi(z)+h\ell_\chi(1)}
 \ll_\nu (\log Q)^2 h z^{\sigma_1-1}
                       +(\log Q) z^2/Q .
\]
The difference of the finite sums is
\(h\sum_{n\le z}\Lambda(n)\chi(n)/n\) plus an error from
\(z<n\le z+h\). Its modulus is at most
\(h(h+1)\log(z+h)/z\). Dividing by \(h\) bounds the tail \(R_\chi(z)\).

\proofstep{Return to the smoothed sum.}

Take \(X=Y_0^J\), \(z=X^{400}\), \(h=z/X^2\) and
\(\nu=\varepsilon/(1600J)\). Multiplication by \(X\) gives
\[
 \abs{X R_\chi(z)}
 \ll_{\varepsilon,J}(\log Q)^2Y_0^{\varepsilon/4}
      +(\log Q)X^{403}/Q+ X^{-1}\log X .
\]
Here \(Y_0^{\varepsilon/6}\le\log Q\le Y_0^{\varepsilon/4}\)
for sufficiently large \(Y_0\). The bound is eventually at most
\(Y_0^\varepsilon\). Also
\(\abs{\ell_\chi(1)}\ll_J\log Y_0\) by the definition of the
tail and the factorial identity \eqref{eq:factorial-mangoldt-bound}.

For \(z=Z\) split the defining sums at \(y\). The identity and the
two bounds, including Proposition~\ref{prop:quadratic-host-uniform}, give
\[
 \begin{aligned}
 G_\chi(y)&=-S_l(y)-yR_\chi(Z),\\
 |G_\chi(y)|&\le \tfrac12 y^{1/2+\tau}+Y_0^\varepsilon
             \le y^{1/2+\tau}\qquad (Y_0\le y\le X).
 \end{aligned}
\]
The last inequality holds for large \(Y_0\) since \(\varepsilon<1/52\). The constants and the choice of the prime apply to
every twist of the bounded list at once. To recover the explicit bound needed when the argument of the weight is shifted, use the estimate with \(0<\tau'<\min(\tau,1/2)\). The centered identity and \eqref{eq:factorial-mangoldt-bound} give, uniformly on the interval,
\[
 |\ell_\chi(1)|\le y^{-1/2+\tau'}+\sum_{n\le y}\Lambda(n)/n
              \le 4\log y
\]
for sufficiently large \(Y_0\).

\end{proof}

The remaining choice is simultaneous in the two primes. Once the first is chosen, the
twists needed for the product must also be included when choosing the second.

\begin{proposition}[A single prime pair for all primitive twists]\label{prop:common-host-prime-pair}
Fix the scale exponents \(J,\varepsilon\) of
Definition~\ref{hyp:host-residue-class}, \(\tau>0\), and a fixed
bounded list of primitive characters. For sufficiently large
\(Y_0\) there exist primes \(q\equiv3\bmod8\) and
\(d\equiv7\bmod16\), with logarithms in \eqref{eq:host-pair-log-windows}, satisfying 
\[
 q\mid d+1,\qquad p\mid d+1\quad(3\le p\le Y_0^{\varepsilon/5},\ p\text{ prime}).
\]
 The bound of
Proposition~\ref{prop:quadratic-host-smoothed} holds for the
primitive twists at \(q,d,qd\) for every \(Y_0\le y\le Y_0^J\).
\end{proposition}

\begin{proof}
\proofstep{Select the pair in order.}

For the first prime choose \(q\equiv3\bmod 8\), with
\(\log q\) in the first interval of \eqref{eq:host-pair-log-windows}.
Keeping this prime, choose \(d\equiv7\bmod16\) with \(d+1\)
divisible by \(q\) and all odd primes at most \(Y_0^{\varepsilon/5}\).
The two logarithm windows are
\begin{equation}\label{eq:host-pair-log-windows}
 \tfrac12Y_0^{\varepsilon/5}<\log q\le Y_0^{\varepsilon/5},\qquad
 \tfrac12Y_0^{\varepsilon/5}(\log Y_0)^2<\log d
           \le Y_0^{\varepsilon/5}(\log Y_0)^2 .
\end{equation}
Put \(Y=Y_0^{\varepsilon/5}\). The conditions on \(d\)
give a reduced class by the Chinese remainder theorem at
\[
 m_0=16q\prod_{3\le p\le Y}p,\qquad
 \log m_0\le (2+o(1))Y.
\]
For the second selection choose \(\log U=Y(\log Y_0)^2\).
The modulus size then gives
\[
 \tfrac14(\log Y_0)^2\log m_0
 \le (\tfrac12+o(1))Y(\log Y_0)^2\le\log U .
\]
The factor \(q\) exceeds \(Z\), and every other prime
divisor of \(m_0\) is at most \(Y\), as required in
Proposition~\ref{prop:quadratic-host-correlation}.

In the second application the base list consists of the original
list and its twists at the already selected \(q\).
If \(F_0\) is the largest original conductor, then the
largest conductor in the second list is at most \(F_0q\).
In the second selection \(F_0\) is fixed and
\[
  \log(F_0q)\le\log F_0+Y\le\tfrac12\log U
\]
for all sufficiently large \(Y_0\). The bad primes for that
list can therefore be discarded by
Proposition~\ref{prop:quadratic-host-smoothed}.
The first choice controls the twists at \(q\), and the
second choice controls the twists at \(d\) and at \(qd\).
This gives one pair for the whole interval \(Y_0\le y\le Y_0^J\).

\end{proof}

\begin{construction}
\begin{tabular}{@{}lll@{}}\toprule
Prime selected & Base list for the selection & New twists controlled\\\midrule
\(q\equiv3\bmod8\) & \(\chi\) & \(\chi\chi_q\)\\
\(d\equiv7\bmod16\) & \(\chi,\ \chi\chi_q\) &
\(\chi\chi_d,\ \chi\chi_q\chi_d\)\\\bottomrule
\end{tabular}
\caption{Selecting one ordered pair for all characters \(\chi\) in the fixed list and all scales \(Y_0\le y\le Y_0^J\). Here \(\chi_p\) is the quadratic character at \(p\), and each product denotes its inducing primitive character.}\label{fig:ordered-prime-twist-list}
\end{construction}

\begin{lemma}[Rescaling the shared interval]\label{lem:host-rescaling}
Let \(q,d\) be the pair of Proposition~\ref{prop:common-host-prime-pair}
and let \([W,W^2]\) be its cohomology interval from
Lemma~\ref{lem:exceptional-twist-window}, with \(Y_0=X^{1/J}\).
For the fixed \(c\ge1\) and \(a_*=c^{1/(2+8\varepsilon)}\), if
\(W\ge\max(a_*^3,a_*^{400})\), the dyadic frame and twist bounds transfer under the change of variables 
\[
 a_*W\le x\le (a_*W)^{3/2},\quad x'=x/a_*\quad\Longrightarrow\quad W\le x'\le W^2.
\]

\end{lemma}

\begin{proof}
\proofstep{Use the same cohomology interval after rescaling.}

Fix the common field factor \(c\ge1\) depending on \(\psi\), and
put \(a_*=c^{1/(2+8\varepsilon)}\). The cohomology window
\([W,W^2]\) will be chosen for this pair, with
\(X^{1/J}\le W\le\sqrt X\) and \(Y_0=X^{1/J}\).
On \([a_*W,(a_*W)^{3/2}]\) write \(x=a_*x'\).
If \(W\ge a_*^3\), then the upper rescaled endpoint is
\((a_*W)^{3/2}\le W^2\).

Since \(a_*\ge1\),
\[
 W\le x'=x/a_*\le x\le W^2,
 \qquad Y_0\le x\le X.
\]
Multiplication by \(c=a_*^{2+8\varepsilon}\) carries the
dyadic frame interval at \(x'\) to the frame interval at
\(x\). Moreover, if \(W\ge a_*^{400}\), then
\(x^{400}=a_*^{400}(x')^{400}\le(x')^{401}\).
In particular every integer \(1\le j<\lfloor x^{400}\rfloor\)
lies in the strict dyadic range \(1\le j<\lfloor(x')^{401}\rfloor\).
The error bound transfers because
\(\exp((x')^\varepsilon/2)\le\exp(x^\varepsilon/2)\).
The same prime pair is used before and after rescaling.
\end{proof}
